%
% Capítulo 7
%
\chapter{Estudo Exploratório com Large Language Models (LLMs) para Interpretação de Disponibilidades} \label{cap:llms}

No âmbito do projeto, foi desenvolvido um estudo com o objetivo de avaliar a viabilidade de utilização de modelos locais e multimodais para interpretar disponibilidades de professores e alunos a partir de diferentes tipos de input.
A motivação para este estudo surgiu da possibilidade de, no futuro, permitir à aplicação receber informação em formatos menos estruturados, como texto livre, imagens, \textit{screenshots}, fotografias manuscritas ou áudio, transformando esses dados em informação útil para o sistema.

A ideia central consistiu em analisar se seria possível utilizar estes modelos como apoio ao registo de disponibilidades, reduzindo a necessidade de inserção exclusivamente manual e estruturada por parte do utilizador.
Até ao momento, o estudo foi dividido em três frentes principais:
\begin{itemize}
    \item extração de texto a partir de imagens;
    \item interpretação de disponibilidades em texto e imagem;
    \item preparação da pipeline necessária para futura análise de áudio.
\end{itemize}

\section{Modelos e Ferramentas Avaliados}

Durante esta fase foram considerados diferentes modelos e ferramentas, com papéis distintos no processamento dos dados:
\begin{itemize}
    \item \textit{llava}~\cite{llava}, utilizado numa fase inicial como modelo multimodal geral;
    \item \textit{glm-ocr:q8\_0}, avaliado como solução local dedicada a OCR;
    \item \textit{qwen2.5vl:3b}~\cite{qwen25-vl}, utilizado como modelo multimodal para interpretação de imagem e texto;
    \item \textit{ffmpeg}, instalado como ferramenta de apoio à preparação de ficheiros áudio;
    \item \textit{Whisper}~\cite{whisper}, considerado para uma futura etapa de transcrição automática de áudio.
\end{itemize}

Entre os modelos testados, o \textit{llava} teve um papel mais exploratório e não foi mantido como solução principal.
O \textit{glm-ocr:q8\_0} foi inicialmente escolhido por ser leve, local e vocacionado para OCR.
Já o \textit{qwen2.5vl:3b} destacou-se como o candidato mais promissor para interpretação de inputs visuais e textuais.

\section{Resultados Obtidos com OCR e Multimodalidade}

Os primeiros testes com o modelo \textit{glm-ocr:q8\_0} mostraram resultados positivos em cenários simples de texto digital limpo.
Em imagens com informação estruturada, o modelo conseguiu ler corretamente nomes e blocos de disponibilidade.
Também apresentou um comportamento aceitável em \textit{screenshots} e texto digital simples.

No entanto, o modelo revelou limitações significativas quando foi necessário estruturar a informação ou lidar com dados menos controlados.
Em tentativas de conversão para JSON, produziu saídas inválidas, alterou informação e traduziu dias da semana para inglês.
Em fotografias manuscritas, o comportamento foi ainda menos robusto, tendo sido observadas alterações de horários, adição de linhas inexistentes e introdução de informação não presente na imagem.
Concluiu-se, por isso, que o \textit{glm-ocr:q8\_0} pode ser útil como OCR simples, mas não constitui uma solução suficientemente fiável para servir de base única à funcionalidade pretendida.

Posteriormente, foram realizados testes com o modelo \textit{qwen2.5vl:3b}, que apresentou um desempenho claramente superior, sobretudo em situações mais próximas de uso real.
Nos testes com fotografias manuscritas, o modelo não se revelou perfeito, mas demonstrou um comportamento mais estável, evitando inventar dias ou horários inexistentes e preservando melhor a estrutura original da informação.
Apesar de alguns erros ortográficos em palavras isoladas, os resultados foram globalmente mais consistentes.

Para além das fotografias manuscritas, o \textit{qwen2.5vl:3b} foi também testado com imagens simples, \textit{screenshots} e texto livre, tendo mostrado capacidade para:
\begin{itemize}
    \item ler corretamente texto digital simples;
    \item interpretar disponibilidades descritas em linguagem natural;
    \item identificar indisponibilidades e preferências;
    \item preservar a informação temporal essencial para posterior utilização no sistema.
\end{itemize}

Com base nos resultados observados, o \textit{qwen2.5vl:3b} constitui, até ao momento, o melhor candidato para suportar uma futura funcionalidade de interpretação de disponibilidades a partir de texto e imagem.

\section{Estado Atual da Componente de Áudio}

A componente de áudio ainda não foi validada de forma completa nesta fase do projeto.
Foi preparado um ficheiro de teste e instalada a ferramenta \textit{ffmpeg} para apoio ao tratamento deste tipo de input.
No entanto, a etapa de transcrição automática ainda não foi concluída, mantendo-se o \textit{Whisper} como a solução mais provável para evolução futura.

Dado que o \textit{qwen2.5vl:3b} não interpreta áudio diretamente, a arquitetura prevista e pensada para este caso passa por uma pipeline em duas etapas:
\begin{itemize}
    \item transcrição do áudio para texto com um modelo de reconhecimento de fala;
    \item interpretação desse texto com o modelo multimodal selecionado.
\end{itemize}

Assim, a abordagem atualmente considerada para este tipo de input corresponde a:
\begin{center}
    \textit{audio} $\rightarrow$ \textit{Whisper} $\rightarrow$ texto transcrito $\rightarrow$ \textit{qwen2.5vl:3b}
\end{center}

\section{Proposta de Integração no Projeto}

Embora esta componente ainda não se encontre integrada na versão beta da aplicação, o estudo realizado já permite delinear uma proposta de arquitetura para evolução futura do sistema.

A solução mais promissora identificada até ao momento assenta nos seguintes fluxos:
\begin{itemize}
    \item texto direto $\rightarrow$ \textit{qwen2.5vl:3b};
    \item imagem simples $\rightarrow$ \textit{qwen2.5vl:3b};
    \item \textit{screenshot} $\rightarrow$ \textit{qwen2.5vl:3b};
    \item fotografia manuscrita $\rightarrow$ \textit{qwen2.5vl:3b} com validação humana;
    \item áudio $\rightarrow$ \textit{Whisper} + \textit{qwen2.5vl:3b}.
\end{itemize}

Esta abordagem permitiria acrescentar à aplicação uma forma mais flexível de recolha de disponibilidades, especialmente útil em contextos em que o utilizador dispõe da informação em formatos menos estruturados.
Ao mesmo tempo, os resultados obtidos mostram que esta integração deverá ser feita com cuidado, prevendo sempre mecanismos de validação e confirmação por parte do utilizador, sobretudo em casos de manuscrito ou áudio.

\section{Síntese do Estudo}

O estudo realizado permitiu concluir que, até ao momento, não foi identificado um único modelo capaz de tratar de forma totalmente fiável todos os tipos de input considerados.
Ainda assim, foi possível identificar uma abordagem tecnicamente promissora para evolução futura da solução,
com destaque para o \textit{qwen2.5vl:3b} como principal candidato para imagem e texto, e para o \textit{Whisper} como etapa complementar no caso do áudio.

Desta forma, esta linha de trabalho não constitui ainda uma funcionalidade concluída da versão beta,
mas representa uma direção de evolução concreta e fundamentada para o projeto, alinhada com o objetivo de tornar o registo de disponibilidades mais natural, acessível e adaptável a diferentes formatos de entrada.
